% kernel/representations.tex

\chapter{Representing the Candidate Without Replacing It}
\label{chap:representations}

The proposed transition is now generated as a continuous curve for
calculation, a sampled polyline for CAD review, and an exchange record for
coordination. These forms make the candidate usable; none becomes the
alignment or owns its identity.

\begin{definition}[Representation]
\label{def:representation}
A representation is a derived description of engineering information
intended for communication, visualization, exchange, computation, or
interaction.
\end{definition}

Each representation must identify what state it expresses, its purpose, the
derivation that produced it, and any approximation. The CAD polyline may be
adequate for visual coordination yet omit the function-family identity and
derivatives required to reconstruct the Berlinish composition. The continuous
calculation model may support curvature evaluation while being unsuitable as
an exchange contract. Fitness is purpose-relative.

Transformation, resampling, or regeneration can therefore produce another
representation of the same candidate state. Conversely, two geometrically
equal files need not refer to the same object. Representation equality is
neither object identity nor engineering acceptance.

An imported historical transition illustrates the other direction. Its file
can remain authoritative source evidence within its scope even when its
sampled geometry is not mathematically optimal. Deriving an editable
continuous model preserves the import and records a new derivation; it does
not retroactively turn the derived model into the original observation.

The candidate is now communicable. To compare it with the railway as found,
the account must next distinguish generated geometry from observed evidence.
